% HRMA User Guide (English) — compile with xelatex
% Build: xelatex HRMA-User-Guide-EN.tex (twice, for the table of contents)
\documentclass[11pt,a4paper]{article}

\usepackage{fontspec}
\usepackage[english]{babel}
\usepackage{geometry}
\geometry{a4paper, margin=2.4cm, top=2.6cm, bottom=2.6cm}
\usepackage{graphicx}
\usepackage{xcolor}
\usepackage{fancyhdr}
\usepackage{titlesec}
\usepackage{enumitem}
\usepackage{tcolorbox}
\usepackage{booktabs}
\usepackage{array}
\usepackage{hyperref}
\usepackage{parskip}
\tcbuselibrary{skins,breakable}

\definecolor{hrmacyan}{HTML}{0891A8}
\definecolor{hrmadark}{HTML}{0A1524}
\definecolor{hrmaink}{HTML}{1B2A3A}
\definecolor{hrmagray}{HTML}{5A6B7A}
\definecolor{hrmapanel}{HTML}{EEF4F8}
\definecolor{hrmawarn}{HTML}{B26A00}

\hypersetup{
  colorlinks=true, linkcolor=hrmacyan, urlcolor=hrmacyan,
  pdftitle={HRMA User Guide}, pdfauthor={Berke Tezgocen}
}

\titleformat{\section}{\Large\bfseries\color{hrmacyan}}{\thesection.}{0.5em}{}
\titleformat{\subsection}{\large\bfseries\color{hrmadark}}{\thesubsection}{0.5em}{}
\titlespacing*{\section}{0pt}{1.4em}{0.6em}

\pagestyle{fancy}
\fancyhf{}
\fancyhead[L]{\small\color{hrmagray}HRMA — Rocket Motor Analysis Suite}
\fancyhead[R]{\small\color{hrmagray}User Guide}
\fancyfoot[C]{\small\color{hrmagray}\thepage}
\renewcommand{\headrulewidth}{0.4pt}
\renewcommand{\footrulewidth}{0pt}

\newtcolorbox{notebox}{
  colback=hrmapanel, colframe=hrmacyan, boxrule=0.6pt, arc=3pt,
  left=8pt, right=8pt, top=6pt, bottom=6pt, breakable}
\newtcolorbox{warnbox}{
  colback=yellow!8, colframe=hrmawarn, boxrule=0.8pt, arc=3pt,
  left=8pt, right=8pt, top=6pt, bottom=6pt, breakable}

\newcommand{\fig}[2]{%
  \begin{center}
  \fbox{\includegraphics[width=0.96\linewidth]{#1}}\\[3pt]
  {\small\color{hrmagray}\itshape #2}
  \end{center}}

\begin{document}

% ================= COVER =================
\begin{titlepage}
\centering
\vspace*{2.5cm}
{\fontsize{54}{60}\selectfont\bfseries\color{hrmadark}HRMA}\\[4pt]
{\Large\color{hrmacyan}Rocket Motor Analysis \& Design Suite}\\[6pt]
{\large\color{hrmagray}Hybrid • Solid • Liquid Rocket Motors}\\[2.2cm]

{\Huge\bfseries\color{hrmadark}User Guide}\\[10pt]
{\large\color{hrmagray}A step-by-step worked example}\\[3cm]

{\large Version 2.6.25 \quad|\quad July 2026}\\[6pt]
{\small\color{hrmagray}Developed by: Berke Tezgocen \quad•\quad Idea \& Testing: Ayberk Cem Aksoy}\\[1.5cm]

\begin{notebox}
\small\color{hrmadark}
\textbf{Scope.} HRMA is a preliminary-design and educational tool built on
closed-form and one-dimensional engineering correlations. It is \emph{not}
a flight-qualification tool. Predicted performance should be cross-checked
against an independent code (CEA / RPA / openMotor) and verified by physical
testing before firing any motor. See ``Scope and Limitations'' at the end of
this guide.
\end{notebox}
\end{titlepage}

% ================= CONTENTS =================
\tableofcontents
\newpage

% ================= 1. INTRODUCTION =================
\section{What is HRMA?}

HRMA (Rocket Motor Analysis Suite) is a desktop application for the
preliminary design and analysis of hybrid, solid, and liquid rocket motors.
It runs a local calculation engine (Flask) inside its own native window and
presents results in a dark-themed, interactive interface: an interactive 3D
model built live from the solver output, an Analysis Deck of engineering
panels, and working CAD / drawing / report exports.

\textbf{Who is it for?} Rocket teams, undergraduate and graduate students,
faculty, and anyone who wants to see motor physics end to end. The program
does not design a motor \emph{from scratch} for you; it produces geometry,
performance and warnings from the parameters you enter, which you then judge
with engineering reasoning.

\begin{notebox}
\small\color{hrmadark}
\textbf{Everything runs locally and offline.} None of your design data
leaves your computer. The only time the program reaches the internet is at
startup, when it asks GitHub whether a newer version exists; if that fails,
the program continues to work normally.
\end{notebox}

\subsection{How to use this guide}

The guide follows a single \textbf{example motor}: a hybrid using HTPB fuel
and liquid N\textsubscript{2}O oxidizer, 1000~N thrust, 10~s burn time. The
same steps apply to solid and liquid motors; the differences are noted in
their sections. All screenshots are taken directly from the application.

% ================= 2. INSTALL =================
\section{Installation and First Launch}

\subsection{Installation}

\begin{itemize}[leftmargin=1.4em]
  \item \textbf{Windows 10/11:} Download \texttt{HRMA-Setup-2.6.25.exe},
  double-click and follow the wizard (Next, Next, Install). The installer is
  per-user; no administrator rights are required. Windows SmartScreen may
  warn because the installer is unsigned: click ``More info'', then ``Run
  anyway''.
  \item \textbf{macOS 11+ (Apple Silicon):} Open
  \texttt{HRMA-Setup-2.6.25-macOS.dmg} and drag \texttt{HRMA} into
  \texttt{Applications}. On first launch, right-click the app and choose
  ``Open'' (Gatekeeper requires this once for unsigned apps).
\end{itemize}

Python and all libraries are bundled; no separate installation is needed.

\subsection{First launch}

The app shows a splash screen within about a second while the calculation
engines load in the background, then opens the home page in its own window.
The home page has three motor-type cards and a link to the formula reference.

\fig{img/en-00-home.png}{The home page: cards to choose a motor type
(Hybrid, Solid, Liquid), with the language selector, Release Notes, User
Guide and Settings links in the top bar.}

\begin{notebox}
\small\color{hrmadark}
\textbf{Language.} The selector in the top-right switches the interface
between English and Turkish; the choice applies to every page, including
chart text. A Turkish edition of this guide ships with the application too.
\end{notebox}

% ================= 3. PAGES =================
\section{Main Pages}

The application has five main pages:

\begin{center}
\begin{tabular}{@{}>{\raggedright}p{3.2cm} p{2.2cm} p{8.5cm}@{}}
\toprule
\textbf{Page} & \textbf{URL} & \textbf{Purpose} \\
\midrule
Home & \texttt{/} & Motor-type selection and formula reference \\
Hybrid Designer & \texttt{/hybrid} & Full hybrid motor design (the most complete page) \\
Solid Designer & \texttt{/solid} & Grain geometry, ballistics, Monte Carlo \\
Liquid Designer & \texttt{/liquid} & Feed system, injector, regen cooling, cycle \\
Formulas & \texttt{/formulas} & The equations used by the solvers (MathJax) \\
\bottomrule
\end{tabular}
\end{center}

All three designer pages share the same structure: an input form, a
Calculate button, a results HUD, standalone panels (transient / injector /
6-DOF), the Analysis Deck, the 3D model, and export buttons.

% ================= 4. WORKED EXAMPLE =================
\section{Worked Example: Designing a Hybrid Motor}

In this section we work through a complete example. Open the \texttt{/hybrid}
page.

\subsection{Step 1 — Enter the parameters}

The form comes with sensible defaults, top to bottom. For our example we use
these values (most are already the defaults):

\begin{center}
\begin{tabular}{@{}p{5cm} p{4cm} p{5cm}@{}}
\toprule
\textbf{Field} & \textbf{Value} & \textbf{Note} \\
\midrule
Motor Name & \texttt{Example-H1} & Used in reports and exports \\
Altitude & \texttt{0 m} & Sea level (ambient pressure auto) \\
Thrust & \texttt{1000 N} & Target thrust level \\
Burn Time & \texttt{10 s} & Total impulse is derived from this \\
O/F Ratio & \texttt{2.5} & Oxidizer/fuel mass ratio \\
Chamber Pressure & \texttt{20 bar} & Combustion chamber pressure \\
Tank Pressure & \texttt{50 bar} & Must exceed chamber pressure enough \\
Fuel & \texttt{HTPB} & Regression coefficients auto-filled \\
Oxidizer & \texttt{N\textsubscript{2}O (liquid)} & Density updates with temperature \\
Injector & \texttt{Showerhead} & Type-specific parameters \\
\bottomrule
\end{tabular}
\end{center}

\fig{img/en-01-form.png}{The input form. The ``Find Optimum'' button sweeps
the O/F ratio for the c*-optimal value. The ``?'' next to each field shows a
short explanation.}

\begin{notebox}
\small\color{hrmadark}
\textbf{0 = automatic.} Enter \texttt{0} into fields such as expansion
ratio, chamber diameter or port diameter and the solver sizes them itself.
Override the defaults with your own data (for example regression
coefficients fitted from a static fire).
\end{notebox}

\subsection{Step 2 — Calculate}

Press \textbf{Calculate}. The solver sends the form to \texttt{/calculate}
and populates the results within a few seconds.

\fig{img/en-02-results.png}{The results HUD (Performance Summary), the
Warnings panel, and the design tables. In our example the Specific Impulse
is 185.6~s, thrust 1000~N, chamber pressure 20~bar, mass flow 0.549~kg/s,
and the minimum structural safety factor 2.89.}

\textbf{Reading the results:}

\begin{itemize}[leftmargin=1.4em]
  \item \textbf{Performance Summary (HUD):} The headline numbers — thrust,
  total impulse, Isp, c*, CF, chamber pressure, throat and exit diameters,
  mass flows, and a safety-state badge.
  \item \textbf{Warnings panel:} Design-criteria messages from the solver
  and the validation system. \textbf{Read these before trusting the
  numbers.} Our example raises three: c* just below its normal range
  ($1325$ m/s), a low expansion ratio ($\varepsilon = 3.4$), and a flash
  boiling risk. These are not errors; they tell you where the design sits at
  a limit and under which assumptions.
  \item \textbf{Design tables:} Full geometry — nozzle contour dimensions
  and angles, grain/port geometry, injector plan, wall thickness.
\end{itemize}

\begin{warnbox}
\small\color{hrmadark}
\textbf{Warnings are a feature, not a defect.} Rather than invent a number
for something it cannot compute, HRMA warns explicitly or says
``not modelled''. When you see a red/orange warning, cross-check that
quantity by hand or with an independent code.
\end{warnbox}

\subsection{Step 3 — Inspect the 3D model}

Below the results is a 3D motor model built live from the geometry the
solver produced. Use the cutaway, dimensions, exploded, heat-map and burn
animation buttons to inspect it.

\fig{img/en-03-3d.png}{The 3D motor simulation (cutaway view). Dimensions
come straight from the solver output: chamber diameter 83.5~mm, throat
21.7~mm, exit 40.2~mm, total length 1377~mm. The timeline at the bottom
animates how the port diameter grows through the burn.}

\begin{notebox}
\small\color{hrmadark}
\textbf{Interactive design mode.} After a calculation, the chamber diameter,
L* and expansion-ratio sliders recompute geometry in about a second; the 3D
model and cross-section update live without a full recalculation.
\end{notebox}

\subsection{Step 4 — Charts and cross-section}

Performance charts (thrust curve, pressure, regression rate and port growth)
are drawn in the app's dark theme, fully offline. The cross-section is a
scaled engineering drawing generated from the same geometry the solver used.

\fig{img/en-05-cross.png}{The motor cross-section with a full legend:
chamber wall, liner, fuel grain, injector plate, nozzle, and the final port
diameter at burnout, all dimensioned in mm.}

\subsection{Step 5 — Export}

The export buttons in the results panel produce: drawing and report PDFs,
CAD files (including STEP cross-section mapping) and CSV/Excel data tables.
Outputs are written to the \texttt{Documents/HRMA} folder.

% ================= 5. OTHER MOTOR TYPES =================
\section{Solid and Liquid Motors}

The solid and liquid pages follow the same pattern; only the inputs are
type-specific.

\begin{itemize}[leftmargin=1.4em]
  \item \textbf{Solid motor (\texttt{/solid}):} Grain geometry (BATES,
  finocyl, slotted, etc.), number of segments, propellant family (KNDX,
  KNSB, APCP and more), ballistics and Monte Carlo uncertainty analysis. The
  burn-rate \texttt{a} and \texttt{n} coefficients come from a database.
  \item \textbf{Liquid motor (\texttt{/liquid}):} Propellant pair, feed
  system type and \textbf{power cycle} (pressure-fed, gas-generator,
  tap-off, staged combustion, FFSC, expander). A cycle power-balance solver
  checks whether the pump and turbine powers close. It also has supercritical
  CH\textsubscript{4}/H\textsubscript{2} regenerative cooling and a
  gas-gas injector path. Combustion thermochemistry is solved with RocketCEA
  at the motor's \emph{real} chamber pressure, O/F ratio and expansion ratio.
\end{itemize}

% ================= 6. SCOPE =================
\section{Scope and Limitations (Please Read)}

HRMA is a capable preliminary-design and educational tool, but knowing its
limits is essential:

\begin{itemize}[leftmargin=1.4em]
  \item \textbf{Preliminary-design level.} The models are closed-form and
  1D correlations. Turbopump maps, cavitation, rotordynamics, bearing life,
  startup/shutdown transients and 2D/3D conjugate heat transfer are
  \emph{not} present.
  \item \textbf{Validation is limited.} Given the real operating point and
  the published geometry, liquid vacuum Isp is reproduced with about
  $\sim1.2\%$ median error over a 14-engine sample. That supports ``predicts
  Isp well at a known point'', \emph{not} ``predicts the flight performance
  of an unknown motor from scratch to that error''. On the hybrid side the
  regression-rate and chamber-pressure errors are much larger (13--35\%).
  For current figures see \texttt{docs/VALIDATION\_STATUS.md}.
  \item \textbf{Not flight qualification.} HRMA alone does not support the
  final design, safety decision or qualification of a motor that will fly.
  That requires fatigue/creep analysis, physical testing and a formal V\&V
  process.
\end{itemize}

\begin{warnbox}
\small\color{hrmadark}
\textbf{The golden rule.} HRMA is a starting point. Cross-check every
critical result with an independent code (CEA / RPA / openMotor) and, where
possible, with a static fire test. Knowing how to use the program is not, by
itself, being a rocket engine engineer.
\end{warnbox}

% ================= 7. UPDATES =================
\section{Automatic Updates}

At startup HRMA checks GitHub Releases for a newer version. If one exists, a
window appears: \textbf{Update now}, \textbf{Later}, or \textbf{Skip this
version}. Choosing ``Update now'' downloads the installer, closes the app,
installs silently, and reopens automatically.

\begin{notebox}
\small\color{hrmadark}
\textbf{New in version 2.6.25.} This release brings three visible
changes. (1) The \emph{Cooling Channels} selection on the hybrid page now
actually enters the calculation; in earlier versions the thermal analysis
always ran uncooled regardless of what you picked. The warning shown when you
select channels states that the coolant film coefficient is taken from
literature and that coolant flow is not verified --- that is a deliberate
limit, not a fault. Chamber material and wall thickness now change the thermal
results as well. (2) Warnings produced by the engines now appear as readable
text in the interface language; previously some of them reached you as raw
codes. (3) The update window shows release notes in your interface language.

\textbf{Update note.} On macOS the update install can take a few
minutes, during which a notification is shown; do not open the app by hand
until it finishes. If your network hits GitHub's request limit (possible on
shared networks), the program now falls back to the ordinary release pages
and still finds the update.
\end{notebox}

\vspace{1cm}
\begin{center}
\color{hrmagray}\small
HRMA v2.6.25 — Rocket Motor Analysis Suite\\
\href{https://github.com/berketez/HRMA}{github.com/berketez/HRMA}
\end{center}

\end{document}
